\documentclass[11pt]{article}
\usepackage[T1]{fontenc}
\usepackage{lmodern,amsmath,amssymb,amsthm,mathtools}
\usepackage[margin=1in]{geometry}
\usepackage{microtype,booktabs,array,longtable}
\usepackage[numbers,sort&compress]{natbib}
\usepackage{xurl}
\usepackage[colorlinks=true,linkcolor=blue,citecolor=blue,urlcolor=blue]{hyperref}
\hypersetup{pdftitle={Verified bounds for positive cubic relaxation gaps}}
\newtheorem{theorem}{Theorem}
\newtheorem{lemma}[theorem]{Lemma}
\newtheorem{proposition}[theorem]{Proposition}
\theoremstyle{remark}
\newtheorem{remark}[theorem]{Remark}
\DeclareMathOperator{\vex}{vex}
\DeclareMathOperator{\cav}{cav}
\DeclareMathOperator{\conv}{conv}
\newcommand{\E}{\mathbb E}
\newcommand{\R}{\mathbb R}
\newcommand{\ind}{\mathbf 1}
\title{Verified bounds for positive cubic relaxation gaps}
\author{}
\date{}
\begin{document}
\maketitle
\begin{abstract}
Luedtke, Namazifar, and Linderoth proved that for bilinear functions
with positive coefficients on the unit cube, the gap of the termwise
relaxation, which uses the exact envelopes of the individual terms, is
less than twice the gap of the convex hull of the graph. We study degree
three.
Let $R(3)$ be the supremum of the ratio of termwise gap to hull gap over
multilinear polynomials of degree at most three with nonnegative
coefficients, all finite nonnegative boxes, and all points with positive
hull gap. We prove
\[
 2.1668\approx\frac{1610000}{743033}\le R(3)\le\frac{31}{12}\approx2.5833.
\]
The upper bound uses one fixed mixture of three probability laws that
preserves every variable mean and bounds every term at once. Its weights
are optimal among mixtures of these three laws for a uniform termwise
guarantee. The lower bound comes from an explicit integer-coefficient
family and a rational Bernstein certificate. Its members lie on unit
cubes and already for $m=36$ have ratio above two, so some positive cubic
polynomials on the unit cube have a larger ratio than every positive
bilinear function there. The exact value of $R(3)$ remains open. An
accompanying Lean project verifies both bounds, the optimality of the
weights, and seven exact finite examples, all stated for the actual graph
hulls.
\end{abstract}

\section{Introduction}\label{sec:intro}
Global solvers for nonconvex problems usually relax each nonlinear term
separately. For a multilinear function, the \emph{termwise} relaxation
replaces each monomial by its own convex and concave envelopes. The
convex hull of the whole graph is stronger, but its standard description
uses all vertices of the box \citep{Luedtke2012}. We measure the loss by
the ratio of the two vertical widths at a point.

\paragraph{Setting.}
Let $B=\prod_{i=1}^n[\ell_i,u_i]$ with $0\le\ell_i\le u_i<\infty$, and let
\[
 f(x)=\sum_{e\in E}a_em_e(x),\qquad m_e(x)=\prod_{i\in e}x_i,
 \qquad a_e\ge0,\quad |e|\le3,
\]
where $E$ consists of distinct supports. Constant and affine terms are
allowed. Write $\vex_B g$ for the greatest convex underestimator of $g$
on $B$, and $\cav_B g$ for its least concave overestimator. The
\emph{hull gap} and the \emph{termwise gap} are
\begin{align*}
 H_B(f;x)&=\cav_B f(x)-\vex_B f(x),\\
 T_B(f;x)&=\sum_{e\in E}\bigl[\cav_B(a_em_e)(x)-\vex_B(a_em_e)(x)\bigr].
\end{align*}
Thus $T_B$ uses the exact hull of each \emph{original} term, and always
$0\le H_B\le T_B$. Define $R(3)$ as the supremum of $T_B/H_B$ over this
class, arbitrary finite dimension, and points with $H_B>0$. Fixed
coordinates, boundary points, and zero coefficients are allowed.

\paragraph{Prior work.}
\citet{Luedtke2012} studied this comparison. For positive coefficients on
a nonnegative box, the concave envelope of $f$ is the sum of the concave
envelopes of its terms, so all of the loss is on the convex side. For
positive bilinear functions on the unit cube, where the termwise and
McCormick relaxations coincide, they proved that $T_B/H_B$ is less than
two, with a factor depending on the chromatic number of the graph of the
terms. They conjectured that for positive multilinear functions on
nonnegative boxes the ratio is bounded by a constant independent of
degree and dimension. Results for one fixed degree, such as ours,
neither confirm nor refute this conjecture, which concerns all degrees at
once. Degree three is the first case beyond their bilinear theorem. Our definitions use exact individual envelopes on each
box; they are not identified with every recursive McCormick formulation,
which can be weaker on a general positive box. We do not claim an
exhaustive literature search for priority.

\paragraph{Results.}
Theorem~\ref{thm:main} gives the two-sided bound on $R(3)$.
\begin{itemize}
\item \emph{Upper bound} (Section~\ref{sec:upper}). A single mixture
 \eqref{eq:mixture} of three mean-preserving laws captures at least
 $12/31$ of every bilinear and cubic term gap at once. This gives
 $T_B\le(31/12)H_B$ on every finite nonnegative box.
\item \emph{Optimality of the weights} (Section~\ref{sec:optimality}).
 No fixed mixture of the same three laws guarantees more than $12/31$
 of every term gap.
\item \emph{Lower bound} (Section~\ref{sec:lower}). The family
 \eqref{eq:family} with positive integer coefficients, together with the
 scalar certificate of Lemma~\ref{lem:minorant}, gives ratios whose
 certified lower bounds converge to $1610000/743033$. Already at $m=36$
 the certified bound without the slack $\delta$ exceeds two, which
 bounds every positive bilinear ratio on the unit cube
 \citep{Luedtke2012}.
\item \emph{Exact finite examples} (Section~\ref{sec:witnesses}) and
 \emph{formal verification} (Section~\ref{sec:formal}).
\end{itemize}
The key tool is the probabilistic description of graph hulls in
Section~\ref{sec:hulls}: on the unit cube, the convex envelope at $x$ is
the minimum of $\E f(X)$ over laws on binary vertices with mean $x$. The
termwise gap lets each term choose its own law, while the hull gap
requires one law for all terms.

\begin{theorem}\label{thm:main}
For every polynomial and point above,
\[
 T_B(f;x)\le\frac{31}{12}H_B(f;x).
\]
Moreover,
\[
 \frac{1610000}{743033}
 =\frac{4830000}{2229099}
 \le R(3)\le\frac{31}{12}.
\]
\end{theorem}

\section{Graph hulls and common laws}\label{sec:hulls}
On the cube, separate affinity gives
\[
 \conv\{(x,f(x)):x\in[0,1]^n\}
 =\conv\{(v,f(v)):v\in\{0,1\}^n\}.
\]
Consequently the lower and upper envelopes at $x$ are the minimum and
maximum of $\E f(X)$ over binary laws with $\E X_i=x_i$. These extrema
are attained. For a nonempty support $e$, its individual endpoints are
\[
 \cav m_e(x)=\min_{i\in e}x_i,
 \qquad
 \vex m_e(x)=\max\left(0,\sum_{i\in e}x_i-|e|+1\right).
\]
The empty support is constant one. The common-threshold law
$X_i=\ind_{\{U\le x_i\}}$, with $U$ uniform on $[0,1]$, attains all
upper endpoints at once. Hence $\cav f=\sum_e a_e\cav m_e$.
For any other law $P$ with these marginals, define its deficiency by
\[
 D_P(e)=\min_{i\in e}x_i-\E_Pm_e(X)\ge0.
\]
If one law has $D_P(e)\ge\alpha H(m_e;x)$ for every nonempty support,
then $H(f;x)\ge\alpha T(f;x)$. The law must be shared by all terms.

\section{The universal upper bound}\label{sec:upper}
Define three global laws with the prescribed marginals:
\begin{itemize}
\item $O$: draw a common uniform $U$ and, independently conditional on
$U$, sample each coordinate with success probability
\[
 p_O(x_i,U)=
 \begin{cases}
 \tfrac12\ind_{\{U\le2x_i\}},&x_i\le1/2,\\
 1-\tfrac12\ind_{\{U\le2(1-x_i)\}},&x_i>1/2.
 \end{cases}
\]
This is the folded orientation law used in the formal construction.
\item $I$: draw independent Bernoulli coordinates of means $x_i$.
\item $B$: call $i$ low when $x_i\le1/2$. A low coordinate is
$X_i=\ind_{\{U\le x_i\}}$. For each high coordinate, independently
conditional on $U$, set $X_i=0$ with probability $1/2$ when
$U\le2(1-x_i)$, and with probability zero otherwise.
\end{itemize}
For $B$, each high coordinate has failure probability
$\tfrac12\,2(1-x_i)=1-x_i$; all three laws preserve all means.
Use the mixture
\begin{equation}\label{eq:mixture}
 P=\frac{18}{31}O+\frac6{31}I+\frac7{31}B.
\end{equation}

\subsection{Cubic terms}
Order the three means as $u\le v\le w$ and put
$T=\min(u,2-v-w)$. Write $D_O,D_I,D_B$ for the deficiencies.
If $v\le1/2$, integrating the conditional product for the first two
coordinates in $O$ gives $\E_O[X_1X_2]=u/2$. Thus $D_O\ge u/2$. Also
$D_I=u(1-vw)\ge u/2$. Here $T=u$, so
$18D_O+6D_I+7D_B\ge12T$.

If $u\le1/2<v$, put $a=1-v$ and $b=1-w$, so $a\ge b\ge0$.
Integration of the conditional success probabilities gives
\begin{equation}\label{eq:onelow}
 D_O=\frac12\min(u,a)+\frac14\min(u,b),\qquad
 D_B=\frac12\min(u,2a)+\frac14\min(u,2b).
\end{equation}
We claim $18D_O+7D_B\ge12\min(u,a+b)$. If $a\le u/2$, the left
side is $16a+8b\ge12(a+b)$. If $b\ge u/2$, it is at least
$18(3u/8)+7(3u/4)=12u$. In the remaining cases $b\le u/2\le a$.
For $a\le u$ the left side is $9a+8b+7u/2$. When $a+b\le u$,
$3a+4b\le7u/2$; when $a+b\ge u$,
$9a+8b\ge8u+a\ge17u/2$. Both give the claim. For $a\ge u$ the
left side is $25u/2+8b\ge12u$. The boundaries are included, and
$D_I\ge0$ completes the mixture estimate.

If $u>1/2$, put $c=1-u$, $a=1-v$, $b=1-w$. Then
$1/2>c\ge a\ge b\ge0$. Under $B$, the failure-union probability
is $c+a/2+b/4$, by integration over the three nested failure intervals.
Direct integration for $O$ gives the same deficiency:
\[
 D_O=D_B=a/2+b/4\ge3(a+b)/8\ge3T/8.
\]
Independence gives the stronger estimate $D_I\ge7T/16$.
Indeed, for $s=a+b\le1$ we have $ab\le s^2/4$ and
$D_I\ge u(s-s^2/4)$. If $s\le u$, then
$D_I/T\ge u(1-u/4)\ge7/16$. If $s\ge u$, then
$D_I/T\ge s-s^2/4\ge u-u^2/4\ge7/16$.
These arguments are needed only for $T>0$; at $T=0$ nonnegativity
suffices. Thus
\[
 18D_O+6D_I+7D_B\ge
 \left(25\frac38+6\frac7{16}\right)T=12T.
\]

\subsection{Bilinear terms and boxes}
For two means $u\le v$, $T=\min(u,1-v)$ and $D_O=T/2$.
If both are low, $D_I\ge T/2$. If just one is low,
$D_B=\min(u,2(1-v))/2\ge T/2$. If both are high,
$D_B=T/2$ and $D_I=uT\ge T/2$. The mixture therefore captures
at least $12T/31$ in every case. Constant and singleton terms have zero
gap. Summing the common-law inequalities proves the cube bound.

For a nonnegative box, substitute $x_i=\ell_i+(u_i-\ell_i)y_i$.
Each original positive monomial expands into nonnegative cube monomials
of degree at most three. The sum of the expanded term widths is at least
the original term width: taking a supremum of a sum is bounded by the
sum of suprema, and taking an infimum has the reverse inequality.
The full hull width is preserved by the affine change of coordinates.
This remains true with zero-width coordinates. Apply the cube result to
the expansion to prove the upper bound in Theorem~\ref{thm:main}.

\section{Optimality of the fixed mixture}\label{sec:optimality}
This section concerns uniform termwise guarantees for fixed mixtures
of $O,I,B$, not the exact value of $R(3)$. Let their nonnegative weights
be $w,r,v$ with sum one. Suppose their mixture captures at least a
fraction $\alpha$ of every bilinear and cubic monomial gap.
Three tests give the following normalized deficiencies:
\begin{center}
\begin{tabular}{lll}
\toprule
Means & Parameter & $(D_O,D_I,D_B)/T$ or its limit\\
\midrule
$(1/2,1/2)$ & fixed & $(1/2,1/2,0)$\\
$(\varepsilon,1-\varepsilon^2,1-\varepsilon^2)$
 & $\varepsilon\downarrow0$ & $(3/8,0,3/4)$\\
$(u,1-u/2,1-u/2)$ & $u\downarrow1/2$, $u>1/2$
 & $(3/8,7/16,3/8)$\\
\bottomrule
\end{tabular}
\end{center}
In the second test take $0<\varepsilon\le1/8$, so $T=2\varepsilon^2$.
The exact independent normalized deficiency is
$\varepsilon-\varepsilon^3/2$. In the third test take $1/2<u\le5/8$,
which keeps all coordinates high and sorted; $T=u$, and the independent
normalized deficiency is $u-u^2/4$. These are tests of the actual
laws, with the threshold equality classified as low in the first test.
It follows that
\[
 \alpha\le\tfrac12-\tfrac v2,\qquad
 \alpha\le\tfrac38+\tfrac{3v}8-\tfrac{3r}8,\qquad
 \alpha\le\tfrac38+\tfrac r{16}.
\]
Average these inequalities with weights $3/31,4/31,24/31$.
The coefficients of $r$ and $v$ cancel, giving $\alpha\le12/31$.
Equation~\eqref{eq:mixture} attains that guarantee. This does not exclude
a better law, coefficient-dependent weights, or an analysis that does
not demand the same deficiency fraction separately for every term.

\section{An analytic family for the lower bound}\label{sec:lower}
Let $m\ge9$ be a multiple of nine. Use disjoint groups $U,V,W$ of
$m$ variables each, with coordinate sums $A,B,C$. Write $E_k(G)$ for
the elementary symmetric polynomial of degree $k$ on a group $G$.
Define
\begin{equation}\label{eq:family}
 f_m=2E_3(W)+3B E_2(W)+2mE_2(V)
       +\frac{5m}{3}AC+\frac{10m}{9}AB+mE_2(U).
\end{equation}
Every included term has degree two or three and a positive integer
coefficient. Set the group means to $1/4,1/2,3/4$, respectively, on the
unit cube $[0,1]^{3m}$.

\subsection{A rational scalar certificate}
Put
\begin{align*}
 F(a,b,c)&=6c^3+27bc^2+18b^2+30ac+20ab+9a^2,\\
 L(a,b,c)&=37a+\frac{79}{2}b+38c-\frac{103}{3},\qquad
 \delta=\frac{901}{120000}.
\end{align*}
\begin{lemma}\label{lem:minorant}
For $(a,b,c)\in[0,1]^3$, $F(a,b,c)\ge L(a,b,c)+\delta$.
\end{lemma}
\begin{proof}
Define
\begin{align*}
 p(c)&=\frac{-972c^4+576c^3+1404c^2-768c+101}{96},\\
 q(c)&=\frac{-78732c^4+212256c^3-203796c^2+82032c-11275}{2976}.
\end{align*}
For $c\le3/10$, let $b_*=13/24-3c^2/4$ and
$g(c)=-49/6+30c-15c^2\le0$. The identity
\[
 F-L-p(c)=9\bigl(a-1+\tfrac{10}{9}(b-b_*)\bigr)^2
       +\tfrac{62}{9}(b-b_*)^2+g(c)(a-1)
\]
proves $F-L\ge p(c)$ when $a\le1$. For every real $a,b,c$,
\begin{align*}
 F-L-q(c)={}&9\left(a+\tfrac{10}{9}b+\tfrac{30c-37}{18}\right)^2\\
 &+\frac{62}{9}\left(b+
 \frac{9(27c^2-79/2)-10(30c-37)}{124}\right)^2.
\end{align*}
It remains to bound $p$ on $[0,3/10]$ and $q$ on $[3/10,1]$.
For each row below, set $t=(c-\ell)/(u-\ell)$. The indicated polynomial
equals $D^{-1}\sum_{i=0}^4v_i\binom4i t^i(1-t)^{4-i}$.
\begin{center}\small
\begin{tabular}{clrl}
\toprule
Polynomial & $[\ell,u]$ & $D$ & $(v_0,v_1,v_2,v_3,v_4)$\\
\midrule
$p$ & $[0,1/5]$ & $60000$ & $(63125,39125,20975,9395,4133)$\\
$p$ & $[1/5,3/10]$ & $240000$ & $(16532,6008,1802,3788,11597)$\\
$q$ & $[3/10,1/2]$ & $7440000$ & $(215357,1285415,1434125,1250375,1008125)$\\
$q$ & $[1/2,3/4]$ & $190464$ & $(25808,18056,7964,9230,15869)$\\
$q$ & $[3/4,1]$ & $190464$ & $(15869,22508,34520,45920,31040)$\\
\bottomrule
\end{tabular}
\end{center}
The Bernstein basis is nonnegative and sums to one. Every $v_i/D$ is
at least $\delta$, proving the lemma by exact polynomial identities.
\end{proof}

\subsection{From counts to graph hulls}
At a binary vertex, put $a=A/m$, $b=B/m$, $c=C/m$. Expansion gives
\begin{equation}\label{eq:expansion}
 \frac{18}{m^3}f_m
 =F(a,b,c)-\frac{18c^2+27bc+18b+9a}{m}+\frac{12c}{m^2}
 \ge F(a,b,c)-\frac{72}{m}.
\end{equation}
Every coupling of the prescribed individual marginals has expected
normalized counts $(1/4,1/2,3/4)$. Taking expectations in
Lemma~\ref{lem:minorant} and \eqref{eq:expansion} therefore yields
\[
 \frac{18}{m^3}\vex f_m\ge\frac{139}{6}+\delta-\frac{72}{m}.
\]
Counting the six support classes and using the individual endpoints gives
\begin{align*}
 \frac{18}{m^3}\cav f_m&=\frac{167}{4}-\frac{153}{4m}+\frac9{m^2},\\
 \frac{18}{m^3}T(f_m)&=\frac{161}{4}-\frac{135}{4m}+\frac6{m^2}.
\end{align*}
Only the $E_3(W)$ terms have nonzero individual lower endpoints.
The actual hull gap is positive, and the preceding inequalities give
\begin{equation}\label{eq:ratio}
 \frac{T(f_m)}{H(f_m)}\ge
 \frac{161/4-135/(4m)+6/m^2}
 {223/12-\delta+135/(4m)+9/m^2}.
\end{equation}
Positivity of the hull gap follows, for example, by comparing the common
upper law with independence: every nonlinear support here has means
strictly between zero and one and a positive coefficient.
As $m$ tends to infinity through multiples of nine, the right side of
\eqref{eq:ratio} tends to
\[
 \frac{161/4}{223/12-901/120000}=\frac{1610000}{743033}.
\]
Every member belongs to the class defining $R(3)$, so this proves the
lower bound in Theorem~\ref{thm:main}. Dropping the slack gives the
simpler bound $483/223$. Already the unrefined finite certificate at
$m=36$ exceeds two. No finite member is claimed to attain either
limiting constant, and no limit of the actual ratios is asserted.

\section{Exact finite examples}\label{sec:witnesses}
The formal project also verifies seven exact finite examples. For the
two-group family, with both groups of size $m$, use
\[
 g_m=A E_2(W)+\frac{5m}{4}E_2(U),
 \qquad x_U=1/2,\quad x_W=3/4.
\]
For the three-group family, use group means $(1/4,1/2,3/4)$ and
coefficients $\gamma=(\gamma_0,\ldots,\gamma_5)$ on the six terms
$(E_3(W),B E_2(W),E_2(V),AC,AB,E_2(U))$.
\begin{center}
\begin{tabular}{lrll}
\toprule
Family & Variables & Coefficients or parameter & Exact $T/H$\\
\midrule
Two groups & 8 & $m=4$ & $27/16$\\
Two groups & 16 & $m=8$ & $21/11$\\
Two groups & 24 & $m=12$ & $99/50$\\
Two groups & 32 & $m=16$ & $135/67$\\
Three groups & 18 & $(2,3,9,10,7,7)$ & $20891/10411$\\
Three groups & 24 & $(2,3,13,12,8,7)$ & $6601/3225$\\
Three groups & 192 & $(2,3,120,105,70,63)$ & $7443345/3445256$\\
\bottomrule
\end{tabular}
\end{center}
Each exact value is certified by an affine lower bound valid on every
integer count vector, a probability law attaining that bound with all
individual marginals preserved, and the common upper law. The last
four ratios exceed two. These finite certificates are independent of
the asymptotic argument above; their full data and kernel proofs are
in the accompanying source distribution.

\section{Formal verification and scope}\label{sec:formal}
The standalone project in \href{formal/README.md}{\texttt{formal/}}
contains the dependency closure of the canonical proofs. Its
\href{formal/COVERAGE.md}{claim map} identifies declarations for the
graph-hull interpretation, the three actual laws, scalar inequalities,
all-box transfer, fixed-mixture obstruction, finite analytic family,
limit and supremum, and finite witnesses. The
\href{formal/VERIFICATION.md}{verification record} gives the pinned Lean
and Mathlib versions, commands, module and declaration counts, axiom
audit, kernel replay, and source fingerprints; source export identity
and the paper build are checked separately. The standard axioms
\texttt{propext}, \texttt{Classical.choice}, and \texttt{Quot.sound} are
allowed; unfinished proofs and additional axioms are rejected.

General coefficient removal, homogeneous approximation of the
supremum, equal-marginal classification, and the general two-group
asymptotic family are separate results outside this paper; the finite
two-group examples above do not depend on them.

\bibliographystyle{plainnat}
\bibliography{references}
\end{document}
